\documentclass{article}
\usepackage{amsmath}
\usepackage{amssymb}
\usepackage{graphicx}
\usepackage{booktabs}
\usepackage{array}
\usepackage{hyperref}
\usepackage{enumitem}
\usepackage{verbatim}

\title{DSC 257R: Unsupervised Learning \\ Homework 10}
\author{}
\date{}

\begin{document}

\maketitle

\section*{Autoencoders}

\begin{enumerate}
\item \textbf{A variational autoencoder for MNIST.} Download the file \texttt{vae.zip} from the course website and unzip it. Look through the Jupyter notebook \texttt{VAE.ipynb}. This fits a variational autoencoder to the MNIST data set and then generates new MNIST-like samples from it. Read through the notebook carefully to get the lay of the land.

The accompanying directory has two pre-trained models in it. One of these has a latent variable dimension of 2, while the other has a latent variable dimension of 8.

\begin{enumerate}
\item By making very minor edits to the notebook, train a VAE with latent dimension 4 and save it. Then generate 100 samples from this VAE and print a large image containing all of them. Note: If the training of the VAE takes too long on your machine, feel free to halve the number of epochs.

\item Also generate a grid of 256 images, generated from latent vectors that are evenly spaced along a random line segment in the 4-dimensional latent space (as in the last section of the notebook). Show two such grids of images.
\end{enumerate}
\end{enumerate}

\section*{Self-supervised Learning}

\begin{enumerate}[resume]
\item \textbf{Word embeddings.} The large number of English words can make language-based applications daunting. To cope with this, it is helpful to have a clustering or embedding of these words, so that words with similar meanings are clustered together, or have embeddings that are close to one another.

John Firth (1957) put it thus:
\begin{quote}
You shall know a word by the company it keeps.
\end{quote}
That is, words that tend to appear in similar contexts are likely to be related. In this mini-project, you will investigate this idea by coming up with an embedding of words that is based on co-occurrence statistics.

The description here assumes you are using Python with NLTK.

\begin{itemize}
\item First, download the Brown corpus (using \texttt{nltk.corpus}). This is a collection of text samples from a wide range of sources, with a total of over a million words. Calling \texttt{brown.words()} returns this text in one long list, which is useful.

\item Remove stopwords and punctuation, make everything lowercase, and count how often each word occurs. Use this to come up with two lists:
\begin{itemize}
\item A vocabulary $V$, consisting of a few thousand (e.g., 5000) of the most commonly-occurring words.
\item A shorter list $C$ of at most 1000 of the most commonly-occurring words, which we shall call context words.
\end{itemize}

\item For each word $w \in V$, and each occurrence of it in the text stream, look at the surrounding window of four words (two before, two after):

\[
w_1 \quad w_2 \quad w \quad w_3 \quad w_4
\]

Keep count of how often context words from $C$ appear in these positions around word $w$. That is, for $w \in V, c \in C$, define

\[
n(w, c) = \# \text{ of times } c \text{ occurs in a window around } w
\]

Using these counts, construct the probability distribution $\operatorname{Pr}(c \mid w)$ of context words around $w$ (for each $w \in V$), as well as the overall distribution $\operatorname{Pr}(c)$ of context words. These are distributions over $C$.

\item Represent each vocabulary item $w$ by a $|C|$-dimensional vector $\Phi(w)$, whose $c$'th coordinate is:

\[
\Phi_c(w) = \max\left(0, \log \frac{\operatorname{Pr}(c \mid w)}{\operatorname{Pr}(c)}\right)
\]

This is known as the (positive) pointwise mutual information, and has been quite successful in work on word embeddings.
\end{itemize}

Now, use this matrix to come up with a 100-dimensional representation of words.

\begin{enumerate}
\item Give a description of your 100-dimensional embedding. The description should be concise and clear, and should make it obvious exactly what steps you took to obtain your word embeddings. Below, we will denote these as $\Psi(w) \in \mathbb{R}^{100}$, for $w \in V$. Also clarify exactly how you selected the vocabulary $V$ and the context words $C$.

\item Investigate your embedding using nearest neighbor. Pick a collection of 20 words $w \in V$. For each $w$, return its nearest neighbor $w' \neq w$ in $V$. A popular distance measure to use for this is cosine distance:

\[
1 - \frac{\Psi(w) \cdot \Psi(w')}{\|\Psi(w)\| \|\Psi(w')\|}
\]

Do the results make any sense?

\item Investigate your embedding by clustering the vocabulary. Using the vectorial representation $\Psi(\cdot)$, cluster the words in $V$ into 100 groups. Clearly specify what algorithm and distance function you use for this, and the reasons for your choices. Look over the resulting 100 clusters. Do any of them seem even moderately coherent? Pick out a few of the best clusters and list the words in them.
\end{enumerate}

Note: The Brown corpus is very small. Current work on word embeddings uses data sets that are many orders of magnitude larger.

\end{document}

